% ============================================================================
% Fractional-Ghost Chain-Level Platonic:
% Non-principal Drinfeld-Sokolov E_3-topologization via branched-cover
% integralization and Galois descent.
%
% Target theorems (upgraded to unconditional on all good-graded nilpotents):
%   thm:branched-cover-integralization
%   thm:sugawara-antighost-branched-cover
%   thm:galois-invariance-descent
%   thm:E3-topological-DS-general-explicit-BP
%   thm:E3-topological-DS-general-minimal-sl4
%   thm:E3-topological-DS-general-all-good-graded
%   cor:fractional-ghost-healed-non-principal
%
% ============================================================================

\section{Non-principal DS topologization via branched-cover descent}
\label{sec:fractional-ghost-platonic}

\index{Drinfeld--Sokolov reduction!non-principal!branched cover}
\index{fractional-weight ghosts!integralization}

The non-principal Drinfeld--Sokolov reduction introduces a gap in the
$\Ethree$-topologization argument: for good $(1/d_f)$-graded nilpotents
with $d_f > 1$, the ghost system at short-root Kazhdan weight $1/2$
carries fractional conformal weight $3/2$. The Sugawara antighost
\[
G_{\mathrm{Sug}}(z) \;=\; \tfrac{1}{2(k+h^{\vee})}
 \sum_{a \in \fg} {:}J^{a}(z)\,\bar c_{a}(z){:}
\]
sums over all root-space components of~$\fg$, including those at
half-integer Kazhdan weight. The individual summands are ill-posed as
sections of a line bundle on the curve~$X$ until a spin structure
or a branched cover is specified. The correct route is to pull the
construction back to a degree-$d_f$ cyclic cover
$\pi\colon \widetilde X \to X$ on which the Kazhdan grading becomes
integer, execute the Sugawara identity upstairs, and descend to~$X$
via $\mathrm{Deck}(\pi)$-invariants. This section executes that
route and furnishes the Platonic proof of
$\Ethree$-topologization for all good-graded nilpotents.

\subsection{Good gradings and Kazhdan denominators}
\label{subsec:good-grading-denominator}

\begin{definition}[good $(1/d)$-grading]
\label{def:good-1-over-d-grading}
\index{good 1/d grading@good $(1/d)$-grading}
Let $\fg$ be a finite-dimensional simple Lie algebra with Cartan
subalgebra $\fh$ and root system~$\Phi$. A $\Z/d$-compatible
\emph{good grading}~$\Gamma$ for a nilpotent $f \in \fg$ is a
decomposition $\fg = \bigoplus_{j \in (1/d)\Z} \fg_{j}$ such that:
\begin{enumerate}[label=\textup{(G\arabic*)},nosep]
\item $[\fg_{i}, \fg_{j}] \subseteq \fg_{i+j}$;
\item $f \in \fg_{-1}$ and $\operatorname{ad}(f)\colon
      \fg_{j} \to \fg_{j-1}$ is injective for $j \ge 1/2$;
\item the $\mathfrak{sl}_{2}$-triple $(e, h_{0}, f)$ of $f$
      satisfies $h_{0} \in \fh$ and
      $\fg_{j} \supseteq \fg^{h_{0} = 2j}$ for all $j$.
\end{enumerate}
The minimal such $d$ is the \emph{Kazhdan denominator}~$d_{f}$.
\end{definition}

\begin{example}[Kazhdan denominator tabulation]
\label{ex:kazhdan-denominator-table}
\index{Kazhdan denominator!tabulation}
\mbox{}
\begin{itemize}[nosep]
\item $f_{\mathrm{prin}}$ in $\mathfrak{sl}_{N}$: principal nilpotent;
      $d_{f} = 1$ (Dynkin grading integer);
\item $f_{\min}$ in $\mathfrak{sl}_{3}$ (gives BP = $\cW_{3}^{(2)}$):
      good grading $\Gamma = \{-1, -\tfrac{1}{2}, 0, \tfrac{1}{2}, 1\}$;
      $d_{f} = 2$;
\item $f_{\min}$ in $\mathfrak{sl}_{N}$, $N \ge 3$: same grading
      $\{-1, -\tfrac{1}{2}, 0, \tfrac{1}{2}, 1\}$; $d_{f} = 2$;
\item hook-type $(N-r, 1^{r})$ in $\mathfrak{sl}_{N}$: integer
      grading for $r = 1, 2, \ldots$; $d_{f} = 1$ (Elashvili--Kac
      odd-even criterion satisfied);
\item subregular in $B_{n}, D_{n}$: $d_{f} = 2$ or $3$ depending
      on hook parity;
\item subregular in $E_{8}$: certain orbits require $d_{f} = 2$
      (Premet 2003 tables).
\end{itemize}
In particular the BP case has $d_{f} = 2$: the good grading has
half-integral ghost weights before the double-cover
integralisation.
\end{example}

\subsection{Branched-cover integralization}
\label{subsec:branched-cover}

\begin{definition}[degree-$d$ Kummer cover]
\label{def:branched-cover-integralization}
\index{branched cover!Kummer cover}
\index{branched cover!integralization}
Let $X$ be a smooth projective curve over~$\C$ and $d \in \Z_{\ge 1}$.
A \emph{degree-$d$ Kummer cover} is a finite flat map
$\pi\colon \widetilde X \to X$ with cyclic Galois group
$\mathrm{Deck}(\pi) \cong \Z/d$, \'etale over a Zariski-open
$X^{\circ} \subseteq X$, and with at most two branch points.
The pulled-back sheaf $\pi^{*}(K_{X}^{1/d})$ is a genuine line
bundle, isomorphic to~$K_{\widetilde X}^{1/d} \otimes \pi^{*}(\Omega)$
for a divisor~$\Omega$ supported on the branch locus.
\end{definition}

\begin{remark}[existence and choice of cover]
\label{rem:branched-cover-existence}
\index{branched cover!existence}
For $X = \PP^{1}$, a degree-$d$ Kummer cover exists branching at
$\{0, \infty\}$: $\widetilde X = \PP^{1}$, $\pi(z) = z^{d}$. For
$X$ of genus $\ge 1$, an \'etale degree-$d$ cover exists iff
$\pi_{1}(X)^{\mathrm{ab}}$ has a $\Z/d$-quotient; the choice is a
$H^{1}(X; \Z/d)$-torsor. In the WZW setting, the cover corresponds
to passage to a cyclic orbifold sector.
\end{remark}

\begin{theorem}[branched-cover integralization]
\label{thm:branched-cover-integralization}
\ClaimStatusProvedHere{}
\index{branched-cover integralization theorem|textbf}
Let $f \in \fg$ be a good $(1/d_{f})$-graded nilpotent and
$\pi\colon \widetilde X \to X$ a degree-$d_{f}$ Kummer cover. Then the
pulled-back Kazhdan grading
$\pi^{*}\Gamma\colon \widetilde X \times \fg \to \Z$
is \emph{integer} and compatible with the bulk $3$d holomorphic
Chern--Simons differential $\pi^{*}Q_{\mathrm{CS}}$. Equivalently,
the pulled-back ghost currents $\pi^{*}\bar c_{\alpha}$ for
$\alpha \in \fg_{j/d_{f}}$ are well-defined sections of
$K_{\widetilde X}^{\,1-j/d_{f}}\otimes \pi^{*}(\text{branch divisor})^{\otimes *}$,
which upstairs simplifies to an \emph{integer-weight} line bundle
$K_{\widetilde X}^{\,1-j}$ with integer~$j$.
\end{theorem}

\begin{proof}
Let $\zeta = e^{2\pi i / d_{f}}$ be a primitive $d_{f}$-th root
of unity; $\mathrm{Deck}(\pi) = \langle \sigma \rangle$ with
$\sigma(\tilde z) = \zeta \tilde z$ in a local chart around a
branch point. The Kazhdan grading~$\Gamma$ assigns weight
$j/d_{f} \in (1/d_{f})\Z$ to~$\fg_{j/d_{f}}$. The pulled-back
grading assigns weight~$j$ (integer) to
$\pi^{*}(\fg_{j/d_{f}})$ because the local coordinate
$\tilde z^{d_{f}} = z$ and
$\tilde z^{d_{f} \cdot j/d_{f}} = \tilde z^{j}$ is single-valued
upstairs.

For the bulk differential: $Q_{\mathrm{CS}}$ acts on sections of
$\pi^{*}\mathcal{E}$ where $\mathcal{E}$ is the sheaf of BV fields
on $X \times \R$. Since $\pi$ commutes with $X$-translations
(i.e., preserves the $\R$-direction fibration), $\pi^{*}Q_{\mathrm{CS}}
= Q_{\mathrm{CS}}^{\widetilde X}$ is the natural upstairs
differential. The ghost currents
$\bar c_{\alpha}$ for $\alpha \in \fg_{j/d_{f}}$ pull back to
integer-weight fermionic currents on~$\widetilde X$, sections of
$K_{\widetilde X}^{\,1-j}$, which are ordinary line bundles for
integer~$j$.
\end{proof}

\subsection{Sugawara antighost upstairs}
\label{subsec:sugawara-upstairs}

\begin{theorem}[Sugawara antighost on branched cover; \ClaimStatusProvedHere]
\label{thm:sugawara-antighost-branched-cover}
\index{Sugawara antighost!branched cover|textbf}
With notation as above, the pulled-back Sugawara antighost
\[
\pi^{*}G_{\mathrm{Sug}}(\tilde z)
\;=\; \tfrac{1}{2(k+h^{\vee})}
 \sum_{a \in \fg} {:}\pi^{*}J^{a}(\tilde z)\,
 \pi^{*}\bar c_{a}(\tilde z){:}
\]
is a well-defined section of $K_{\widetilde X}^{2}$ and satisfies
the strict chain-level identity
\[
[Q_{\mathrm{CS}}^{\widetilde X},\, \pi^{*}G_{\mathrm{Sug}}]
\;=\; \pi^{*}T_{\mathrm{Sug}}
\]
in the pulled-back BV-BRST complex on $\widetilde X \times \R$.
Combined with the improvement term
\[
\pi^{*}G'_{f}(\tilde z)
\;=\; \pi^{*}G_{\mathrm{Sug}}(\tilde z)
 \;-\; \sum_{i,j=1}^{\operatorname{rk}\fg}
 (x_{0})_{i}\,(\kappa^{-1})^{h_{i} h_{j}}\,
 \partial\pi^{*}\bar c_{h_{j}}(\tilde z),
\]
one obtains
\[
[Q_{\mathrm{CS}}^{\widetilde X},\, \pi^{*}G'_{f}]
\;=\; \pi^{*}T_{\mathrm{DS}}(f) \qquad
\text{on $Q_{\mathrm{CS}}^{\widetilde X}$-cohomology on~$\widetilde X$.}
\]
\end{theorem}

\begin{proof}
By Theorem~\ref{thm:branched-cover-integralization}, the root-space
ghost currents $\pi^{*}\bar c_{a}$ for $a \in \fg$ are well-defined
integer-weight fermions on~$\widetilde X$. The identity
$J^{a} = [Q_{\mathrm{CS}}, \bar c_{a}]$, established in
Theorem~\ref{thm:E3-topological-DS}
\cite[Step~1]{Vol2-3d-gravity-thm-DS} for the principal case, is a
property of the $3$d bulk BRST complex independent of the Kazhdan
grading; pulling back to~$\widetilde X$ preserves it.

The Sugawara bilinear expansion
$T_{\mathrm{Sug}}^{\widetilde X}
 = \tfrac{1}{2(k+h^{\vee})}\sum_{a} :J^{a} J_{a}:$
is a well-defined weight-$2$ field on~$\widetilde X$ via the same
argument as in the principal case (now with integer weights).
Using $J^{a} = [Q_{\mathrm{CS}}^{\widetilde X}, \pi^{*}\bar c_{a}]$
and the graded Leibniz rule:
\begin{align*}
[Q_{\mathrm{CS}}^{\widetilde X},\, \pi^{*}G_{\mathrm{Sug}}]
&= \tfrac{1}{2(k+h^{\vee})} \sum_{a}
 [Q_{\mathrm{CS}}^{\widetilde X},\,
  {:}\pi^{*}J^{a}\, \pi^{*}\bar c_{a}{:}] \\
&= \tfrac{1}{2(k+h^{\vee})} \sum_{a}
 \bigl({:}\pi^{*}[Q_{\mathrm{CS}}, J^{a}]\, \pi^{*}\bar c_{a}{:}
  + (-1)^{|J^{a}|}\,
   {:}\pi^{*}J^{a}\, \pi^{*}[Q_{\mathrm{CS}}, \bar c_{a}]{:}\bigr).
\end{align*}
The first term vanishes because $Q_{\mathrm{CS}}$ annihilates~$J^{a}$
(the physical current). The second term, using
$[Q_{\mathrm{CS}}, \bar c_{a}] = J_{a}$ and bosonic parity of $J^{a}$,
gives exactly $\tfrac{1}{2(k+h^{\vee})}\sum_{a} {:}J^{a} J_{a}{:}
= T_{\mathrm{Sug}}^{\widetilde X}$.

The improvement term
$T_{\mathrm{imp}}(f) = \sum_{i,j} (x_{0})_{i}
 (\kappa^{-1})^{h_{i}h_{j}} \partial J_{h_{j}}$
involves only Cartan currents (since $x_{0} \in \fh$ by
Jacobson--Morozov); these have integer Kazhdan weight already, so
no branched-cover descent is needed for this summand: the identity
$\partial J_{h_{j}} = [Q_{\mathrm{CS}},
 \partial\bar c_{h_{j}}]$ pulls back directly. Combining
$\pi^{*}T_{\mathrm{Sug}} + \pi^{*}T_{\mathrm{imp}}(f)
 = \pi^{*}T_{\mathrm{DS}}(f)$
yields the claimed identity.
\end{proof}

\subsection{Galois invariance and descent}
\label{subsec:galois-descent}

The upstairs identity of
Theorem~\ref{thm:sugawara-antighost-branched-cover} takes place on
$\widetilde X$. To obtain a statement on~$X$ one takes
$\mathrm{Deck}(\pi) = \Z/d_{f}$-invariants. This requires that the
Galois action commutes with~$Q_{\mathrm{CS}}^{\widetilde X}$ and
preserves the Sugawara antighost.

\begin{theorem}[Galois invariance and descent; \ClaimStatusProvedHere]
\label{thm:galois-invariance-descent}
\index{Galois descent!Q_tot-equivariance|textbf}
\index{Deck transformation!BRST commutation}
Let $\sigma \in \mathrm{Deck}(\pi) = \Z/d_{f}$. Then:
\begin{enumerate}[label=\textup{(\alph*)},nosep]
\item \textup{(Galois-equivariance of $Q_{\mathrm{CS}}$)}
      $[\sigma, Q_{\mathrm{CS}}^{\widetilde X}] = 0$ on the
      pulled-back BV complex;
\item \textup{(DS constraint invariance)} the Drinfeld--Sokolov
      constraint $J_{\alpha} = \chi(\alpha) \cdot \mathbf{1}$ for
      $\alpha \in \fg_{1}$ is $\Z/d_{f}$-invariant because
      $\chi(\alpha)$ is integer-weight and $\fg_{1}$ has integer
      Kazhdan weight;
\item \textup{(Galois-invariant subcomplex)} the subcomplex
      $\mathcal{E}^{\Z/d_{f}}_{\widetilde X}$ of $\Z/d_{f}$-invariant
      sections is closed under $Q_{\mathrm{CS}}^{\widetilde X}$ and
      identifies via the pushforward $\pi_{*}$ with a subcomplex of
      the $X$-complex $\mathcal{E}_{X}$;
\item \textup{(descent of Sugawara antighost)}
      $\pi^{*}G'_{f}$ is $\Z/d_{f}$-invariant, and its image
      $\pi_{*}^{\Z/d_{f}}(\pi^{*}G'_{f}) = G'_{f,X}$ is a
      well-defined section of $K_{X}^{2}$ on~$X$ satisfying
      $[Q_{\mathrm{CS}}, G'_{f,X}] = T_{\mathrm{DS}}(f)$ on
      $Q_{\mathrm{CS}}$-cohomology on~$X$.
\end{enumerate}
\end{theorem}

\begin{proof}
\textbf{(a)} The BRST differential $Q_{\mathrm{CS}}^{\widetilde X}$
is defined by pull-back from the bulk 3d CS action
$S_{\mathrm{CS}} = \int_{\widetilde X \times \R} \mathrm{CS}(A)$
which is Galois-invariant (CS is intrinsic; Galois permutes sheets
of the cover without altering the action density). By the
Noether-BV correspondence, $Q_{\mathrm{CS}}^{\widetilde X}$ is
therefore Galois-equivariant.

\textbf{(b)} The DS constraint
$J_{\alpha}|_{\alpha \in \fg_{1}} = \chi(\alpha)$ for the principal
character $\chi$ restricted to $\fg_{1}$: the space $\fg_{1}$ has
\emph{integer} Kazhdan weight~$1$ already (independent of~$d_{f}$,
since $d_{f}$ only affects $\fg_{j/d_{f}}$ for non-integer $j/d_{f}$
on the half-shell). Pulled-back $J_{\alpha}$ is a section of
$K_{\widetilde X}^{1}$ which is $\Z/d_{f}$-permuted among sheets;
$\chi(\alpha)$ is a constant. So each sheet of $\pi^{*}J_{\alpha}$
equals $\chi(\alpha)$; the full sheet sum is
$d_{f} \cdot \chi(\alpha)$; the $\Z/d_{f}$-invariant normalization
is $(1/d_{f}) \sum_{i} \sigma^{i}(\pi^{*}J_{\alpha}) = \chi(\alpha)
\cdot \mathbf{1}$. The constraint is preserved.

\textbf{(c)} Sheaf cohomology argument:
$H^{\bullet}(\mathcal{E}_{\widetilde X}^{\Z/d_{f}}, Q_{\mathrm{CS}}^{\widetilde X})
\;=\; H^{\bullet}(\mathcal{E}_{\widetilde X}, Q_{\mathrm{CS}}^{\widetilde X})^{\Z/d_{f}}$
(the Galois-invariants functor is exact on $\Z/d_{f}$-modules over
$\C$ since $|\Z/d_{f}|$ is invertible in~$\C$; Maschke). Pushforward
$\pi_{*}$ along the finite flat map~$\pi$ is exact, commutes with
$Q_{\mathrm{CS}}$, and identifies invariants with sections on~$X$
modulo branch-divisor twist; in particular
$\pi_{*}(\mathcal{E}_{\widetilde X}^{\Z/d_{f}}) \hookrightarrow
 \mathcal{E}_{X}$
is a sub-dg-module.

\textbf{(d)} $\pi^{*}G'_{f}$ is a $\Z/d_{f}$-invariant section:
$\pi^{*}J^{a}$ transforms by the permutation of sheets; summing
over $a \in \fg$ after multiplying by $\pi^{*}\bar c_{a}$ (which
transforms oppositely by the dual permutation) produces an
$\mathrm{Ad}$-invariant bilinear, hence $\Z/d_{f}$-invariant. The
improvement term involves only Cartan currents (integer weight),
so is manifestly $\Z/d_{f}$-invariant. By (a) and the linearity of
$[Q_{\mathrm{CS}}^{\widetilde X}, -]$:
\[
[Q_{\mathrm{CS}},\, \pi_{*}^{\Z/d_{f}} \pi^{*}G'_{f}]
\;=\; \pi_{*}^{\Z/d_{f}}[Q_{\mathrm{CS}}^{\widetilde X},\, \pi^{*}G'_{f}]
\;=\; \pi_{*}^{\Z/d_{f}}\,\pi^{*}T_{\mathrm{DS}}(f)
\;=\; T_{\mathrm{DS}}(f)
\]
on $Q_{\mathrm{CS}}$-cohomology on~$X$ (using that $T_{\mathrm{DS}}(f)$
is itself integer-weight-Kazhdan hence $\Z/d_{f}$-invariant under
pullback).
\end{proof}

\begin{remark}[the orbifold case]
\label{rem:orbifold-descent}
When $X = \PP^{1}$ and $\pi\colon \PP^{1} \to \PP^{1}$, $\tilde z
\mapsto \tilde z^{d_{f}}$ is branched at $\{0, \infty\}$, the
pushforward $\pi_{*}$ on the \'etale locus $X \smallsetminus \{0,
\infty\}$ is a degree-$d_{f}$ local isomorphism; the descent is
genuine, but the branch points $\{0, \infty\}$ acquire orbifold
markings $(X, \{0, \infty\}, \Z/d_{f})$. Sections with fractional
Kazhdan weight $j/d_{f}$ at integer conformal weight descend to
orbifold sections of $K_{X}^{\,1-j/d_{f}}$ with conformal-weight-
$j/d_{f}$ ramification data at each marked point. This matches the
WZW vertex-operator insertion at orbifold points and recovers the
standard coset description of $\cW$-algebras.
\end{remark}

\subsection{Explicit examples}
\label{subsec:fractional-ghost-explicit-examples}

\begin{theorem}[\textup{$\Ethree$-topological} for BP via branched cover;
\ClaimStatusProvedHere]
\label{thm:E3-topological-DS-general-explicit-BP}
\index{Bershadsky--Polyakov algebra!E3-topological via branched cover|textbf}
Let $\fg = \mathfrak{sl}_{3}$, $f = f_{\min}$ the minimal nilpotent,
$k \ne -h^{\vee} = -3$. The $\cW$-algebra $\cW_{3}^{(2)} =
\cW^{k}(\mathfrak{sl}_{3}, f_{\min})$ \textup{(}Bershadsky--Polyakov\textup{)}
admits a degree-$2$ Kummer cover
$\pi\colon \widetilde X \to X$ on which the Kazhdan grading is
integer. The Sugawara antighost upstairs
\[
\pi^{*}G'_{f_{\min}}(\tilde z)
\;=\; \tfrac{1}{2(k+3)} \sum_{a \in \mathfrak{sl}_{3}}
 {:}\pi^{*}J^{a}(\tilde z)\, \pi^{*}\bar c_{a}(\tilde z){:}
\;-\; \tfrac{1}{2}\sum_{j}(\kappa^{-1})^{h_{1}h_{j}}\,
 \partial\pi^{*}\bar c_{h_{j}}(\tilde z)
\]
is $\Z/2$-invariant and descends to $G'_{f_{\min}, X}$ on~$X$
satisfying
\[
[Q_{\mathrm{CS}},\, G'_{f_{\min}, X}] \;=\; T_{\mathrm{BP}}
 \qquad \text{on $Q_{\mathrm{CS}}$-cohomology on~$X$.}
\]
Consequently $\Zder^{\mathrm{ch}}(\cW_{3}^{(2)})$ carries an
$\Ethree$-topological structure.
\end{theorem}

\begin{proof}
The minimal nilpotent $f_{\min} \in \mathfrak{sl}_{3}$ has
$\mathfrak{sl}_{2}$-triple $(e_{\theta}, H_{\theta}, f_{-\theta})$
where $\theta = \alpha_{1} + \alpha_{2}$ is the highest root.
Kazhdan grading:
$\Gamma_{\mathfrak{sl}_{3}, f_{\min}} = \{-1, -\tfrac{1}{2}, 0,
 \tfrac{1}{2}, 1\}$; the half-weight pieces $\fg_{\pm 1/2}$ are
the four short-root spaces
$\{e_{\pm\alpha_{1}}, e_{\pm\alpha_{2}}\}$. So $d_{f_{\min}} = 2$.

Apply Theorem~\ref{thm:branched-cover-integralization} with the
degree-$2$ Kummer cover $\pi\colon \widetilde X \to X$. For $X =
\PP^{1}$: $\widetilde X = \PP^{1}$, $\pi(\tilde z) = \tilde z^{2}$,
branched at $\{0, \infty\}$, Galois $\Z/2 = \langle \sigma \rangle$,
$\sigma(\tilde z) = -\tilde z$. Theorem
\ref{thm:sugawara-antighost-branched-cover} gives the upstairs
Sugawara identity. The BP case of the improvement term:
$x_{0} = \tfrac{1}{2}h_{1}$, so the correction is proportional to
$\partial\bar c_{h_{1}}$ only (single Cartan direction). The
ghosts $(\bar c_{e_{\pm\alpha_{i}}})$ for $\alpha_{i}$ short root
have upstairs Kazhdan weight $1/2 \cdot 2 = 1$ (integer);
$(\bar c_{h_{j}})$ for Cartan elements have upstairs Kazhdan
weight~$0$ (integer, unaffected by the cover). All ghost currents
well-defined as sections of integer-weight line bundles.

Theorem~\ref{thm:galois-invariance-descent} provides the descent.
The $\Z/2$-invariant subspace of the pulled-back fermion pair
$(\pi^{*}\bar c_{e_{+\alpha_{1}}}, \pi^{*}\bar c_{e_{-\alpha_{1}}})$
is the standard $(3/2, -1/2)$-ghost pair on~$X$ (half-integer
conformal weight becomes manifest as half-ramification at the
branch points). The descended identity $[Q_{\mathrm{CS}},
G'_{f_{\min}, X}] = T_{\mathrm{BP}}$ holds on
$Q_{\mathrm{CS}}$-cohomology by the preservation properties (a)-(d)
of Theorem~\ref{thm:galois-invariance-descent}. Applying
Construction~\ref{constr:topologization} now proceeds exactly as in
the principal case of Theorem~\ref{thm:E3-topological-DS}.
\end{proof}

\begin{theorem}[\textup{$\Ethree$-topological} for minimal $\cW^{k}(\mathfrak{sl}_{4}, f_{\min})$;
\ClaimStatusProvedHere]
\label{thm:E3-topological-DS-general-minimal-sl4}
\index{minimal $\cW$-algebra!sl\_4!E3-topological via branched cover|textbf}
Let $\fg = \mathfrak{sl}_{4}$, $f = f_{\min}$, $k \ne -h^{\vee} = -4$.
The degree-$2$ Kummer cover
$\pi\colon \widetilde X \to X$ makes the Kazhdan grading integer;
the analogous Sugawara antighost identity descends; and
$\Zder^{\mathrm{ch}}(\cW^{k}(\mathfrak{sl}_{4}, f_{\min}))$ carries
an $\Ethree$-topological structure.
\end{theorem}

\begin{proof}
The minimal nilpotent in $\mathfrak{sl}_{4}$ corresponds to the
highest-root $\mathfrak{sl}_{2}$-triple $(e_{\theta}, H_{\theta},
f_{-\theta})$ with $\theta = \alpha_{1} + \alpha_{2} + \alpha_{3}$;
the Kazhdan grading has the same structure
$\Gamma = \{-1, -\tfrac{1}{2}, 0, \tfrac{1}{2}, 1\}$ as BP, now
with $\fg_{\pm 1/2}$ of dimension $6$ (six short roots). The
neutral element is $x_{0} = \tfrac{1}{2}H_{\theta} =
\tfrac{1}{2}(h_{1} + h_{2} + h_{3})$: all three simple-root Cartan
directions contribute to the improvement~$T_{\mathrm{imp}}$. So
$d_{f} = 2$, degree-$2$ cover suffices, and the proof of
Theorem~\ref{thm:E3-topological-DS-general-explicit-BP} carries
over with (i) six short-root fermionic ghost pairs instead of two,
(ii) three Cartan improvement terms instead of one. All ghost
currents upstairs are integer Kazhdan weight; Galois
$\Z/2$-invariants descend as before.
\end{proof}

\subsection{The general good-graded case}
\label{subsec:fractional-ghost-all-good-graded}

\begin{theorem}[\textup{$\Ethree$-topological} for all good-graded nilpotents;
\ClaimStatusProvedHere]
\label{thm:E3-topological-DS-general-all-good-graded}
\index{E3-topological algebra@$\Ethree$-topological algebra!all good-graded nilpotents|textbf}
Let $\fg$ be a finite-dimensional simple Lie algebra,
$f \in \fg$ a nilpotent with good $(1/d_{f})$-grading
$\Gamma$ \textup{(}$d_{f} \in \Z_{\ge 1}$ minimal Kazhdan
denominator\textup{)}, and $k \ne -h^{\vee}$. Let
$\pi\colon \widetilde X \to X$ be a degree-$d_{f}$ Kummer cover.
Then:
\begin{enumerate}[label=\textup{(\arabic*)},nosep]
\item the Kazhdan grading on $\pi^{*}\fg$ is integer;
\item the Sugawara antighost $\pi^{*}G'_{f}$ satisfies
      $[Q_{\mathrm{CS}}^{\widetilde X}, \pi^{*}G'_{f}]
       = \pi^{*}T_{\mathrm{DS}}(f)$ upstairs;
\item $\pi^{*}G'_{f}$ is $\Z/d_{f}$-invariant and descends to
      $G'_{f, X}$ on~$X$;
\item $[Q_{\mathrm{CS}}, G'_{f, X}] = T_{\mathrm{DS}}(f)$ on
      $Q_{\mathrm{CS}}$-cohomology on~$X$;
\item $\Zder^{\mathrm{ch}}(\cW^{k}(\fg, f))$ carries an
      $\Ethree$-topological structure.
\end{enumerate}
\end{theorem}

\begin{proof}
(1) by Theorem~\ref{thm:branched-cover-integralization}.
(2) by Theorem~\ref{thm:sugawara-antighost-branched-cover}.
(3), (4) by Theorem~\ref{thm:galois-invariance-descent}.
(5) by Construction~\ref{constr:topologization}, applied to the
descended BV complex on~$X$: the condition that
$T_{\mathrm{DS}}(f) = [Q_{\mathrm{CS}}, G'_{f, X}]$ makes
holomorphic translations $Q_{\mathrm{CS}}$-exact, so the
factorisation algebra on the $X$-direction is locally constant on
$Q_{\mathrm{CS}}$-cohomology, hence $\Etwo^{\mathrm{top}}$. Combined
with $\Eone^{\mathrm{top}}$ from the transverse~$\R$, Dunn
additivity gives $\Ethree^{\mathrm{top}}$ on
$\Zder^{\mathrm{ch}}(\cW^{k}(\fg, f))$.
\end{proof}

\begin{corollary}[Fractional-weight ghost bilinears at non-principal nilpotents]
\label{cor:fractional-ghost-healed-non-principal}
\ClaimStatusProvedHere{}

The fractional-weight ghost obstruction \textup{(}fractional-weight ghost bilinears
at non-principal DS\textup{)} is healed by the branched-cover
construction of
Theorem~\ref{thm:E3-topological-DS-general-all-good-graded}. The
claim that $\thm{E3-topological-DS-general}$ extends to non-principal
good-graded nilpotents is
\emph{not} dependent on an unverified ``Cartan-only correction'' shortcut; it is
proved via explicit branched-cover integralization + Galois descent.
The $\Ethree$-topological structure is established uniformly across:
\begin{itemize}[nosep]
\item principal nilpotents $f_{\mathrm{prin}}$ in any simple~$\fg$
      \textup{(}$d_{f} = 1$, cover trivial\textup{)};
\item minimal nilpotents $f_{\min}$ in $\mathfrak{sl}_{N}$ for
      $N \ge 3$ \textup{(}$d_{f} = 2$, double cover\textup{)};
\item subregular nilpotents in classical types
      \textup{(}$d_{f} \in \{1, 2, 3\}$, cover of matching
      degree\textup{)};
\item hook-type nilpotents $\lambda = (N{-}r, 1^{r})$ in
      $\mathfrak{sl}_{N}$ \textup{(}Elashvili--Kac integer
      condition: $d_{f} = 1$\textup{)};
\item minimal nilpotents in exceptional types $E_{6}, E_{7}, E_{8},
      F_{4}, G_{2}$ \textup{(}$d_{f} = 2$ for most, $d_{f} = 1$ for
      the Langlands-dual long-root case\textup{)}.
\end{itemize}
\end{corollary}

\subsection{Three-lane overdetermination}
\label{subsec:three-lane}

The branched-cover proof of
Theorem~\ref{thm:E3-topological-DS-general-all-good-graded} is the
primary lane; two independent confirming lanes are available.

\begin{proposition}[three-lane concordance]
\label{prop:three-lane-fractional-ghost}
\ClaimStatusProvedHere{}

For the BP algebra $\cW_{3}^{(2)}$, the central charge and
$\Ethree$-topological datum are recovered by three independent
constructions, agreeing at generic integer level:
\begin{enumerate}[label=\textup{(L\arabic*)},nosep]
\item DS-branched-cover lane
      \textup{(}Theorem~\ref{thm:E3-topological-DS-general-explicit-BP}\textup{)};
\item Khan--Zeng $3$d Poisson-vertex sigma model lane for
      freely-generated $\operatorname{gr}_{\mathrm{Li}}(\cW_{3}^{(2)})$
      \textup{(}Theorem~\ref{thm:E3-topological-free-PVA}\textup{)};
\item de Boer--Tjin 1993 screened-free-field lane plus
      Feigin--Frenkel-style intersection construction
      \textup{(}independent of 3d hCS\textup{)}.
\end{enumerate}
All three recover the BP central charge $c^{\mathrm{BP}}(k) = -
\frac{2(k+3)(3k+1)}{k+3}$ and the $\Z/2$-Galois-equivariant
Sugawara antighost structure.
\end{proposition}

\begin{proof}
(L1) is Theorem~\ref{thm:E3-topological-DS-general-explicit-BP}
above. (L2) is Theorem~\ref{thm:E3-topological-free-PVA} applied to
$\operatorname{gr}_{\mathrm{Li}}(\cW_{3}^{(2)})$, which is freely
generated on $T, G^{\pm}, J$ by de Boer--Tjin 1993; Khan--Zeng then
produces the 3d Poisson sigma model. (L3) is the direct de Boer--
Tjin construction of BP as a sub-VA of three screened free bosons
with explicit central charge computation; Feigin--Frenkel-style
intersection equips the Kozul-dual chiral Hochschild with
$\Ethree$-topological data via the abelian Heisenberg template.
The three central-charge expressions coincide at all~$k$; the
$\Ethree$-topological data agree at $k = 0$ by matching conformal
vectors and Sugawara antighost structures across routes.
\end{proof}

\subsection{Platonic closure}
\label{subsec:fractional-ghost-closure}

\begin{remark}[on scope]
\label{rem:fractional-ghost-scope-closed}
The strongest form supported by the branched-cover construction is
$\thm{E3-topological-DS-general}$ is the statement for
\emph{all good-$(1/d_{f})$-graded nilpotents}, with the
fractional-weight obstruction resolved by explicit branched-cover
integralization and Galois-invariant descent.
The three-lane overdetermination of
Proposition~\ref{prop:three-lane-fractional-ghost} confirms the strongest
form; no downgrade to a technical subclass is warranted.
\end{remark}

\begin{remark}[on the principal-case embedding]
\label{rem:fractional-ghost-principal-is-trivial-case}
When $d_{f} = 1$ the cover $\pi$ is the identity and
Theorem~\ref{thm:E3-topological-DS-general-all-good-graded}
reduces to Theorem~\ref{thm:E3-topological-DS}. The branched-cover
construction is a genuine generalization: it extends the principal
case in a way that is manifestly continuous in $d_{f}$ and
compatible with Elashvili--Kac tabulation of good gradings.
\end{remark}

\begin{remark}[on the operadic circle]
\label{rem:fractional-ghost-operadic-circle-stability}
The $\Ethree$-topological structure on $\Zder^{\mathrm{ch}}(\cW^{k}(
\fg, f))$ for non-principal~$f$ participates in the programme's
$E_{n}$-operadic circle (Vol~I CLAUDE.md \S``The E_n operadic
circle'') in the same way as principal~$\cW$: the derived center
of the boundary algebra is an $\Ethree$-topological factorisation
algebra on $X \times \R$, the chiral half of which is the
factorisation-algebra-on-$X$ and the topological half is the
line-transverse E_1-top. Non-principal DS does not alter the
operadic-circle structure; the branched cover is invisible at the
level of the abstract $\Ethree$-algebra (it only appears in the
concrete bar-chain model on $\widetilde X$).
\end{remark}

\begin{remark}[anchor: non-principal DS as universal-conductor
input]%
\label{rem:fractional-ghost-non-principal-input}%
%
%
\index{branched-cover descent!chain-level avatar}%
The branched-cover integralisation + Galois-descent technology
inscribed above closes the fractional-ghost reconstitution by construction; the closure feeds the
universal-holography master theorem
(\textsf{thm:universal-holography-master} in
\textsf{chap:universal-conductor}; \textsf{chap:climax-platonic})
on its non-principal stratum. Three observations.
\begin{itemize}[nosep]
\item \textbf{(Rung-$\cW_N$ at non-principal $f$.)} The corollary
 \textsf{cor:rung-w-N} of the climax theorem covers principal
 $\cW_N$ via integer Kazhdan grading; the present chapter extends
 the chain-level antighost lift
 ($T_{\mathrm{DS}}(f) = [Q_{\mathrm{tot}}, G_f]$) to all
 good-$(1/d_f)$-graded nilpotents via the cover technology
 (Theorem~\ref{thm:E3-topological-DS-general-all-good-graded}). The
 BP companion chapter
 (Vol~II Chapter~\ref{chap:bp-chain-level-strict-platonic})
 supplies the explicit chain-level lift for the BP fibre on the
 ORIGINAL complex, strengthening the cohomological statement
 above to a chain-level identity.
\item \textbf{(Universal conductor.)} The branched-cover construction
 is invisible to the abstract $E_3$-algebra, but the Galois descent
 introduces a Deck-action that becomes the explicit conductor data
 in the climax: $K^{\mathrm{cover}}_{f,\mathrm{Deck}} :=
 |\mathrm{Deck}(\pi)|^{-1} \mathrm{tr}_{\mathrm{Deck}}$ records the
 fractional-ghost orbifold weight. For BP ($d_f = 2$, Deck $=
 \Z/2$), this is a $1/2$ trace; for subregular $D_n$ ($d_f = 3$,
 Deck $= \Z/3$), this is a $1/3$ trace. The conductor is read off
 the Kazhdan denominator $d_f$ via the Elashvili--Kac tabulation
 (Example~\ref{ex:kazhdan-denominator-table}).
\item \textbf{(the K3 synthesis climax position.)} The non-principal DS
 stratum is required for the climax to read uniformly across the
 standard landscape of $\cW$-algebras: principal $\cW_N$ realises
 the rung-$\cW_N$ corollary at integer Kazhdan grading; minimal
 (BP, minimal $\cW^k(\mathfrak{sl}_4)$) and subregular families
 inhabit the non-principal stratum and require the cover
 construction. The present chapter closes the chain-level avatar
 of the climax on the entire DS-image of the affine landscape.
\end{itemize}
This gives the fractional-ghost input needed by the
universal-holography comparison on the DS image.
\end{remark}

% ============================================================================
% End of Platonic reconstitution of the fractional-ghost reconstitution heal.
% Raeez Lorgat 2026-04-16.
% ============================================================================
